\subsection{The source's uniform regularity induction}
\label{subsec:ftb-uniform-regularity-induction}

This subsection supplies the uniform algebraic degree bound used
in the actual torsion-form construction.
It is the source's induction on the number of independent linear
coordinates; no Hilbert--Samuel replacement is used.

For \(A=K[X_i:i\in\Sigma]\), with \(\Sigma\) finite, put
\[
 m(e)=\dim_KA_e=\#\{\alpha\in\mathbb N^\Sigma:|\alpha|=e\},
 \quad H_I(e)=\dim_K I_e,\quad
             \langle I_e\rangle=(F:F\in I_e)\subset A.
\]
A homogeneous ideal is \emph{saturated here} if
\[
                    (X_iF\in I\text{ for every }i)
                                      \Longrightarrow F\in I.
 \tag{R7}\label{eq:ftb-regularity-saturation}
\]
Iterating this condition also removes any power of the
irrelevant ideal.
The Hilbert function \(H_I\) is that of the ideal, rather
than the homogeneous quotient.

\begin{lemma}[Uniform regularity in the exact source form]
\label{lem:ftb-uniform-regularity}
For every finite \(\Sigma\) and polynomial
\(Q\in\mathbb Q[T]\), there is \(d_0\), before all fields
\(K\), such that every homogeneous ideal \(I\subset K[X_i]\)
satisfying \eqref{eq:ftb-regularity-saturation} and
\(H_I(e)=Q(e)\) for all sufficiently large \(e\) has
\[
 \begin{split}
 H_I(e)&=Q(e)\quad(e\ge d_0),\\
 I_e&\subset\langle I_{d_0}\rangle\quad(e\ge d_0).
 \end{split}
 \tag{R8}\label{eq:ftb-source-regularity}
\]
The bound is uniform in the characteristic and in \(I\).
\end{lemma}
\begin{proof}
We first work over an infinite field and allow an arbitrary
rational-valued eventual function \(q(e)\), instead of
\(Q(e)\).
Induct on a number \(k\) satisfying
\[
                              m(1)\le H_I(1)+k.
\]
We construct \(D_k(q)\), before \(K,I\), such that for
\(e\ge D_k(q)\),
\[
 H_I(e)=q(e),\qquad
                  I_{e+1}\subset\langle I_e\rangle .
 \tag{R9}\label{eq:ftb-regularity-one-step}
\]
For \(k=0\), the inequality makes \(I_1=A_1\).
Saturation applied to \(1\) gives \(I=A\).
Choose \(D_0(q)\) so that \(m(e)=q(e)\) thereafter
whenever this eventual equality is true; if it is not,
the case is empty and take \(D_0(q)=0\).
All forms of degree \(e+1\) are generated by forms of
degree \(e\), giving the second assertion.

We record the two actual homogeneous-ideal operations
used in the induction.
First, a proper saturated homogeneous ideal over an
infinite field has a degree-one nonzerodivisor \(h\)
on \(A/I\).
The zero divisors are the union of its finitely many
associated primes.
If every linear form were a zero divisor, finite
subspace avoidance would put all linear forms in one
associated prime.
Write that prime as the annihilator of a nonzero
class \(F\bmod I\).
Then every \(X_iF\in I\); saturation gives \(F\in I\),
a contradiction.
Choose \(h\) outside that finite union.

Second, for every homogeneous \(J\) there is a homogeneous
\(J'\supset J\) satisfying
\eqref{eq:ftb-regularity-saturation} and agreeing with
\(J\) in all sufficiently high degrees.
Here is the source's Noetherian construction.
Choose a maximal homogeneous overideal with that
high-degree agreement.
If \(L=(J':(X_i)_i)\), then \(J'\subset L\), \(L\) is
homogeneous, and \(L/J'\) contributes only finitely
many degrees.
To see the last point, choose finitely many generators
of \(L\).
Since every \(X_iF\in J'\) for \(F\in L\), the positive
degree part of any polynomial multiplier times \(F\)
lies in \(J'\).
Modulo \(J'\), only constant multiples of these finitely
many generators remain.
Taking a degree larger than all their degrees gives
zero.
Maximality therefore gives \(L=J'\), as required.

Suppose \(k\ge1\) and \(I\ne A\).
Choose \(h\) as above, let \(J=I+(h)\), and take \(J'\)
by the second construction.
Since \(h\notin I\), one has
\(H_{J'}(1)>H_I(1)\), so \(J'\) meets the induction
inequality with \(k-1\).
Multiplication by \(h\) gives the exact dimension identity
\[
                  H_J(e+1)+H_I(e)=H_I(e+1)+m(e).
 \tag{R10}\label{eq:ftb-hyperplane-dimension}
\]
Indeed \(J_{e+1}=I_{e+1}+hA_e\), their intersection
is \(hI_e\), and multiplication by \(h\) is injective.

The eventual function of \(J'\) is therefore
\[
                         q'(e)=q(e)+m(e-1)-q(e-1).
\]
Put
\[
 d_1=D_{k-1}(q'),\quad d_2=D_{k-1}(q),\quad
 B=\left\lceil m(d_1)-q(d_1)\right\rceil_+,
 \qquad
 D_k(q)=\max(d_2,d_1+B+1).
 \tag{R11}\label{eq:ftb-regularity-recursion}
\]
The \(d_2\) term covers the possible case \(I=A\).
For the proper case, define
\[
       g(e)=H_I(e)-q(e),\qquad
       \delta(e)=H_{J'}(e)-H_J(e)\in\mathbb N.
\]
For \(e\ge d_1\), induction on \(J'\) and
\eqref{eq:ftb-hyperplane-dimension} give
\[
                       g(e+1)=g(e)-\delta(e+1).
 \tag{R12}\label{eq:ftb-regularity-defect-recurrence}
\]
The sequence \(g\) is eventually zero, and
\(g(d_1)\le B\).
It follows from the recurrence that \(g(e)\ge0\) for
\(e\ge d_1\): compare with a later index where it is zero.

Also, for \(e\ge d_1\), if \(\delta(e)=0\), then
\(J'_e=J_e\).
Induction gives
\(J'_{e+1}\subset\langle J'_e\rangle
 =\langle J_e\rangle\subset J\);
hence \(\delta(e+1)=0\).
Thus a nonzero defect at \(d_1+r\) would force a decrease
of at least \(1\) at each of the preceding \(r\) steps.
It would imply \(0\le g(d_1+r)\le B-r\).
At \(r=B+1\) this is impossible.
Consequently \(\delta(e)=0\) for every \(e\ge d_1+B+1\).
The recurrence then makes \(g\) constant there; its
eventual value zero makes it zero there.

It remains to lift generation back through \(h\).
For \(e\) in this range and \(F\in I_{e+1}\), induction
puts \(F\in\langle J'_e\rangle=\langle J_e\rangle\).
Since \(J=I+(h)\), write
\(F=A+hC\), with \(A\in\langle I_e\rangle\).
As \(F,A\in I\), regularity of \(h\) gives \(C\in I\).
Take homogeneous parts of degree \(e+1\):
\[
                           F=A_{e+1}+hC_e.
\]
The first term remains in the homogeneous ideal
\(\langle I_e\rangle\), and the second belongs to it as
well, since \(C_e\in I_e\).
This proves \eqref{eq:ftb-regularity-one-step} with
the exact recursion \eqref{eq:ftb-regularity-recursion}.

Take \(k=m(1)\), which works for every ideal.
Then iterate the one-step generation to obtain
\eqref{eq:ftb-source-regularity} over infinite fields.

Finally remove the infinite-field hypothesis by the
source's coefficient-contraction argument.
Let \(L\) be an algebraic closure of \(K\).
For a \(K\)-linear functional \(\lambda:L\to K\), apply
\(\lambda\) coefficientwise to polynomials, denoting this
by \(c_\lambda\).
Define
\[
 I_L=\{F\in L[X_i]:c_\lambda(F)\in I
                              \text{ for every }\lambda\}.
\]
This is an ideal.
Closure under multiplication by a coefficient \(a\in L\)
uses \(c_\lambda(aF)=c_{\lambda\circ(a\,\cdot)}(F)\);
closure under \(X_i\) is immediate.
The degree pieces are the scalar extensions of \(I_e\):
using a basis of the finite \(K\)-span of a polynomial's
coefficients and its dual coefficient functionals proves
both inclusions.
Their dimensions are therefore unchanged.
Homogeneity and saturation also transfer, since
contraction commutes with homogeneous parts and with
multiplication by \(X_i\).
The infinite-field result applies to \(I_L\) with the
same \(d_0\).

To descend its generation, contraction maps
\(\langle(I_L)_e\rangle\) into
the coefficientwise extension of \(\langle I_e\rangle\).
This follows by expanding polynomial multipliers into
coefficients and monomials and using the two
contraction identities just stated.
Choose \(\lambda\) with \(\lambda(1)=1\).
Applied to a polynomial over \(K\), its contraction is
the same polynomial.
Thus membership obtained over \(L\) descends to
membership in \(\langle I_e\rangle\).
Dimensions and the eventual function already agree.
This proves the lemma for every field.
\end{proof}

\paragraph{Source concordance and exact uniformity.}
Every step above is the mechanism in
\texttt{src/Development/AbelianVarietyTheory/Proofs/ModuliProofs/%
ExistsUniformRegularityBound.lean}
(\texttt{d8f397f6bd3e041101d2001dddcedc9ff05ba5f5}).
The proof uses its literal
\(\mathrm{piece}\), \(\mathrm{genIn}\),
\(\mathrm{exists\_regular\_form}\),
\(\mathrm{exists\_saturation}\),
\(\mathrm{numeric}\), \(\mathrm{bound\_infinite}\), and
coefficient-contraction base change.
It chooses \(d_0\) before all fields and ideals for the
given finite coordinate set and polynomial.
It supplies no host-dependent or window-dependent
choice.  Its application still requires the actual
theta embedding to have the displayed saturated ideal
and the exact eventual Hilbert polynomial; those joins
are separate parts of the torsion-form producer.
